\section{Comparison with Earlier Work}\label{sec:earlier-work}

\subsection{Earlier Work}

Let $H$ consist of the edges of a nonempty strong clique in $G$ and
their endpoints. Write $\Delta=\Delta(G)$ and $h=\Delta(H)$. Fix an
edge $xy$ of $H$. Every edge of $H$ meets $X=N_G(x)\cup N_G(y)$.
Indeed, $x,y\in X$, and any edge disjoint from $xy$ must have an
endpoint adjacent to $x$ or $y$. Thus $|E(H)|\leq |X|h\leq2\Delta h$.
Faudree, Gy\'arf\'as, Schelp and Tuza~\cite{FGST1990} asked what
would happen if this count came arbitrarily close to $2\Delta^2$.
The set $X$ would have close to $2\Delta$ vertices, and most of them
would have close to $\Delta$ edges of $H$ leading outside $X$, since
an edge inside $X$ is counted at both ends. Consider two such edges
$zu,z'u'$ with distinct $z,z'\in X$ and distinct $u,u'\notin X$.
If none of $zz',zu',z'u$ is present, then $uu'$ must be present.
The authors count these required edges at the outside neighbors of
a fixed $z$, together with the edges from those neighbors back to
$X$. The two counts exceed the available degree when
$|E(H)|>(2-\varepsilon)\Delta^2$ for sufficiently small absolute
$\varepsilon>0$. Hence $|E(H)|\leq(2-\varepsilon)\Delta^2$.

For simple graphs, \'{S}leszy\'{n}ska-Nowak~\cite{SleszynskaNowak2016}
used all the edges at a vertex $v$ with $d_H(v)=h$. At most $\Delta h$
edges of $H$ meet $N_G(v)$. Let $F$ consist of the remaining edges.
For every $a\in N_H(v)$, the set $N_G(a)\cap V(F)$ covers $E(F)$,
since each edge of $F$ must be joined to $va$ and neither endpoint is
adjacent to $v$. A vertex belonging to $j$ of these covers already
has $j$ neighbors outside $V(F)$ and hence has degree at most
$\Delta-j$ in $F$. When $h\geq\frac{3}{4}\Delta$, her count using
these covers gives $|E(F)|\leq\Delta^2-\frac{1}{2}\Delta h$.
Adding the edges meeting $N_G(v)$ gives
$|E(H)|\leq\Delta^2+\frac{1}{2}\Delta h\leq\frac{3}{2}\Delta^2$.
For smaller $h$, the bound $|E(H)|\leq2\Delta h$ already suffices.

Faron and Postle~\cite{FaronPostle2019} kept track of the degree sum
at the two ends of each edge of $H$. Write $\sigma=\sigma_G(H)$ for
the largest of these sums, with both degrees measured in $G$.
They chose $x$ of maximum degree in $H$ and, among its neighbors
in $H$, chose $y$ of largest degree in $H$. Their double counting
argument starts by summing the degrees in $H$ of the vertices in
$N_G(x)\cup N_G(y)$ and subtracting the edges counted twice.
For the second count, consider the edges from outside this union to
vertices adjacent to $x$ but not $y$, and those to vertices adjacent
to $y$ but not $x$. Each set forms a
smaller bipartite subgraph of $H$. In the first subgraph neither
endpoint of an edge is adjacent to $y$. For that edge to have a
joining edge to every edge of $H$ at $y$, each neighbor of $y$ in
$H$ must be adjacent to one of its endpoints. Those neighbors lie outside
the subgraph, so restricting $G$ to the vertices of this subgraph
lowers the endpoint degree sum by at least $d_H(y)$. The second
subgraph gives the corresponding decrease of at least $d_H(x)$.

Combining the two counts, Faron and Postle showed that an estimate
with coefficient $\beta\in[\frac{1}{4},\frac{1}{3}]$ for smaller
bipartite subgraphs of $H$ gives the estimate
$|E(H)|\leq\frac{1+\beta}{4}\sigma^2$ for general $H$.
Here the bipartite estimate is $|E(H')|\leq\beta\sigma_{G'}(H')^2$,
where $H'$ is the smaller bipartite subgraph and $G'=G[V(H')]$.
The edges of $H'$ still form a strong clique in $G'$ because all
their endpoints are retained. At $\beta=\frac{1}{3}$, the coefficient
$\frac{1+\beta}{4}$ is again $\frac{1}{3}$, so induction on $|E(H)|$
gives $|E(H)|\leq\frac{1}{3}\sigma^2\leq\frac{4}{3}\Delta^2$.

Kumar, Mohar and Pragada~\cite{KumarMoharPragada2026} improved the
bipartite estimate in this reduction. Suppose that $H$ is bipartite.
Choose $v$ with $d_H(v)=h$, put $A=N_H(v)$ and let
$B=V(H)\setminus N_G[v]$. For an edge $xy$ of $H[B]$, every
$a\in A$ must be adjacent to $x$ or $y$, since $va$ and $xy$ need
a joining edge. Thus $d_{G[B]}(x)+d_{G[B]}(y)\leq\sigma-h$,
and induction bounds $|E(H[B])|$ by $\beta(\sigma-h)^2$.
The same edges from $A$ to $B$ also enter the count at $A$.
Those outside $H$ use degree at $A$ without contributing to $|E(H)|$.
The authors subtract these edges when counting the edges of $H$
meeting $N_G[v]$. Combining this count with the induction bound
reduces the estimate to an inequality in $h/\sigma$, which they
verify for $\beta=\frac{620}{1987}$. The reduction above and
$\sigma\leq2\Delta$ then give
$|E(H)|\leq(1+\beta)\Delta^2=\frac{2607}{1987}\Delta^2$.

A different question is when a five-cycle must occur and what its
vertices tell us about the rest of the graph. Chung, Gy\'arf\'as,
Tuza and Trotter~\cite{CGTT1990} considered simple graphs in which
the entire edge set is a strong clique. After deleting isolated
vertices, such a graph has no induced subgraph consisting of two
disjoint edges. In the triangle-free, nonbipartite case, a shortest
odd cycle must have length five, since a longer induced cycle contains
two edges with no joining edge. Every vertex has exactly two neighbors
on the five-cycle. Vertices with the same two neighbors form one of
five independent sets. Every edge between successive sets is present,
so if their sizes in cyclic
order are $a_0,\ldots,a_4$, then the edge count is
$\sum_{i=0}^4 a_i a_{i+1}$ and $a_{i-1}+a_{i+1}\leq\Delta$ for
each $i$, with indices modulo five. Maximizing this sum gives
$f(\Delta)$. In the bipartite case each part contains a vertex
adjacent to the whole other part, so there are at most $\Delta^2$
edges. If there is a triangle, they find a maximum clique such that
every vertex outside it has a neighbor in it. Grouping the outside
vertices according to their neighbors in this clique and counting
the edges gives fewer than $f(\Delta)$ edges.

Cames van Batenburg, Kang and Pirot~\cite{CamesKangPirot2020}
allowed $H$ to be a proper submultigraph and assumed that $G$ has
no five-cycle. Again take $v$ with $d_H(v)=h$. Consider an edge
$xy$ of $H$ with both ends outside $N_G[v]$. Each $a\in N_H(v)$
must be adjacent to $x$ or $y$. If distinct $a,b\in N_H(v)$
were adjacent to $x$ and $y$, respectively, then $vaxybv$ would
be a five-cycle. Thus, when $|N_H(v)|\geq2$, one endpoint of $xy$
is adjacent to every vertex of $N_H(v)$. The authors count the
edges at these endpoints, then subtract the degree already used by
their neighbors in $N_H(v)$ when counting the remaining clique edges.
Together with the case $|N_H(v)|=1$, this gives
$|E(H)|\leq h(\sigma-h)\leq\frac{1}{4}\sigma^2\leq\Delta^2$.
Consequently a larger strong clique forces some five-cycle in $G$.
For our argument we need an induced five-cycle in the simple support
of $G$ with an edge of $H$ on each of its five pairs of endpoints.

\subsection{Highlights of the Proof}\label{subsec:method-highlights}

The main idea is to use the structure of a counterexample at odd
maximum degree to derive an inequality on edge multiplicities that
contradicts the weighted flag inequalities. Suppose that $G$ has odd
maximum degree $D$ and contains a strong clique of $m=f(D)+1$ marked edges.
For $D\geq9$, \textbf{Lemma~\ref{lem:counterexample-cycle}} gives an
induced five-cycle in the simple support of $G$. Each edge of the cycle
has the same endpoints as a marked edge of $G$. This allows us to
use the five sets suggested by the blow-ups of $C_5$ in the work
of Chung, Gy\'arf\'as, Tuza and Trotter~\cite{CGTT1990}, even when
the host has unmarked edges and parallel edges.

For each such cycle $C$, let $U_0,\ldots,U_4$ be the sets defined
in \textbf{Subsection~\ref{subsec:odd-sketch}}. All edges between
$U_i$ and $U_{i+1}$ are present in the simple support, with indices
taken modulo five. Thus at the $i$th cycle vertex at least
$|U_{i-1}|+|U_{i+1}|$ of its incident edges are already accounted for.
Summing over the five vertices leaves at most
$R(C)=5D-2\sum_{i=0}^4|U_i|$ further incidences.
When $D\geq11$, counting the marked edges in and outside these
sets shows that $R(C)\leq\frac{D+1}{2}$ would give $m\leq f(D)$.
Consequently $R(C)>\frac{D+1}{2}$ on every such cycle of a
counterexample. This use of restrictions on small induced subgraphs
to sharpen a flag bound is related to the stability arguments of
Pikhurko, Slia\v{c}an and Tyros~\cite{PikhurkoSliacanTyros2019}.

Let $s$ be the number of distinct pairs of vertices joined by
marked edges. Summing the lower bound on $R(C)$ over the cycles
and substituting it into the flag inequalities for the simple
support gives an upper bound on $s$. For example, for odd
$D\geq15$ we obtain
\[
  \frac{s}{D^2}\leq 1.221800627
  <\frac{11}{9}\leq\frac{m}{D^2}.
\]
Once one marked edge has been counted on each of the $s$ pairs,
the other $m-s$ marked edges are parallel copies. An upper bound
on $s$ is therefore a lower bound on the fraction $\frac{m-s}{m}$
of such copies. At $D=15$, for instance, the degree estimates give
$R(C)\geq11$ on every induced red five-cycle. Using this bound in the
inequalities for the simple support, and counting the extra copies at
their endpoints, gives $0.007\leq\frac{m-s}{m}\leq0.365$
in \textbf{Section~\ref{sec:exact}}.
With these restrictions, the weighted flag inequalities in the
$D=15$ row of \textbf{Table~\ref{tab:multi-odd}} give
$\frac{m}{D^2}<0.827$. Yet $m=f(15)+1=275$ requires
$\frac{m}{D^2}=\frac{11}{9}$. The other rows give the contradiction
for the remaining odd $D\geq9$.


